
\begin{comment}
    
    


\noindent\textbf{Regression Tests for Transplantation.}~
\Mohsen{ME!}

\noindent\textbf{Regression Tests for Fuzzing.}~
\Mohsen{GIGI!}

\noindent\textbf{WebAssembly Fuzzing.}~
Research on WASM fuzzing has explored diverse approaches. Structure and rule-guided methods, such as WasmCFuzz \cite{zhang2024wasmcfuzz}, RGFuzz \cite{park2025rgfuzz}, and SwFuzz \cite{wang2024s}, synthesize valid Wasm programs to explore deep compiler and runtime behaviors. Mutation-based tools, including Wapplique \cite{zhao2024wapplique}, wasm-mutate \cite{arteaga2022wasm}, WasmFuzzer \cite{jiang2022wasmfuzzer}, FreeWavm \cite{qian2025freewavm}, and Waltzz \cite{zhang2025waltzz}, transform existing modules to diversify test inputs and uncover vulnerabilities. Static and program-analysis approaches, such as WasmChecker \cite{zhuang2024wasmchecker} and WBSan \cite{wu2025wbsan}, detect type errors, memory issues, and undefined behavior without relying solely on fuzzing, while formal verification with WasmRef-Isabelle \cite{watt2023wasmref} provides a mechanized interpreter as a trusted oracle. Differential testing for Wasm has also matured: Wadiff \cite{zhou2023wadiff} compares engine behavior on the same inputs to expose inconsistencies, and Wasmaker \cite{cao2024wasmaker} generates semantics-aware binaries for cross-engine checking. More recently, LLM-assisted systems such as LWDIFF \cite{zhou2025lwdiff} and DrWASI \cite{zhang2025drwasi} guide differential testing and triage across Wasm runtimes and WASI implementations. Complementary studies include Random Testing of WASI \cite{stanley2024random} and the SoK “Analysis Techniques for WebAssembly” \cite{harnes2024sok}, which establish baselines and contextualize security challenges. Our work is most closely related to mutation-based fuzzing \cite{zhao2024wapplique,arteaga2022wasm,jiang2022wasmfuzzer,qian2025freewavm,zhang2025waltzz} but differs by deriving seeds directly from real-world bug reports and using them to guide fuzzing and multi-runtime execution, bridging practical vulnerability sources with systematic exploration.

\vspace{0.5ex}
\noindent\textbf{Seed Corpora Mining and Analysis.}~
Prior work shows that mining real defects, curating provenance, and shaping the corpus significantly influence fuzzing outcomes \cite{just2014defects4j,widyasari2020bugsinpy,tomassi2019bugswarm,herrera2021seed,pailoor2018moonshine,lemieux2018fairfuzz,manes2021survey}. Curated bug datasets such as Defects4J, BugsInPy, Bugs.jar, and BEARS package real bug–fix pairs with fault-revealing tests for reproducible use \cite{just2014defects4j,widyasari2020bugsinpy,saha2018bugsjar,madeiral2019bears}. \tname{} proposes a new seed corpus built from bug histories and regression tests rather than a new benchmark. CI-mined collections like BugSwarm extend this idea to large-scale fail-pass harvesting and document reproducibility risks such as environment drift and flaky builds \cite{tomassi2019bugswarm}. To control these risks, we re-run the fault-revealing tests on the cited versions to reproduce failures before inclusion \cite{tomassi2019bugswarm}. Empirical work on seed selection shows that initial seeds significantly affect outcomes and recommends publishing exact seeds and evaluating across corpora \cite{herrera2021seed,chen2019enfuzz}. \tname{} follows this guidance by comparing seed corpora under fixed configurations. We will release the \tname{} corpus to support reuse, as EnFuzz highlights the value of sharing effective seeds across fuzzers \cite{herrera2021seed,chen2019enfuzz}. Prior work also reduces redundancy through corpus distillation (MoonShine) and coverage-guided prioritization (FairFuzz) \cite{pailoor2018moonshine,lemieux2018fairfuzz,manes2021survey}. In contrast, we apply deduplication so that failing inputs remain small and non-overlapping. For evaluation, Magma centers metrics on ground-truth bugs rather than coverage alone; we mirror this where ground truth exists, even though our seeds originate from historical bugs rather than injections \cite{hazimeh2020magma}. 
%
Structure-aware methods (\eg{}, Skyfire~\cite{wang2017skyfire}, LearnFuzz~\cite{godefroid2017learnfuzz}, and Superion~\cite{wang2019superion}) show that grammar-guided and learned features improves fuzzing on structured inputs.
\tname{} real bug seeds are intended to complement such generators as starting corpora.

\vspace{0.5ex}
\noindent\textbf{Evaluation in Fuzzing}~
Comparative studies show that protocol design, statistical treatment, seed control, and benchmark choice significantly influence measured fuzzer performance \cite{klees2018evaluating}, \cite{arcuri2011statpract}, \cite{bohme2022reliability}. Following this guidance, \tname{} uses repeated trials and controlled randomness, five independent 24-hour runs per configuration and publishes exact seeds to enable replication \cite{klees2018evaluating}, \cite{arcuri2011statpract}, \cite{herrera2021seed}. Shared platforms such as FuzzBench demonstrate reproducible head-to-head campaigns \cite{metzman2021fuzzbench}, while GreenBench studies shorter runs and ranking stability \cite{ounjai2023greenbench}. We follow the reproducibility conventions but use fixed 24-hour campaigns in all experiments. Bug-centric suites such as Magma and LAVA provide ground truth beyond coverage. Accordingly, where ground truth exists we report bug outcomes in addition to coverage \cite{hazimeh2020magma}, \cite{dolangavitt2016lava}. Mutation-analysis work reports low mutant-kill rates after separating coverage and mutation phases. Therefore, \tname{} maintains detailed failure records for each seed, including provenance and failure metadata, to support richer oracles and analysis \cite{gorz2023mutation}. Mutator-level studies show that operator choices can change outcomes and to avoid confounding effects, \tname{} uses default mutator settings from existing fuzzers across all trials \cite{kukucka2024mutator}. Metrics-driven frameworks such as UNIFUZZ encourage multi-dimensional reporting, we adopt this perspective by presenting both bug and coverage metrics under fixed configurations and shared seeds \cite{li2021unifuzz}.

\vspace{0.5ex}
\noindent\textbf{Transplantation of Tests.}~
Test transplantation moves tests from a donor to a recipient by adapting harnesses, aligning configuration, and preserving or adjusting oracles, and has been shown to be effective for JavaScript engines, database systems, mobile apps. \Mar{Confused. Test transplantation used for things that are not tests?}, ~\cite{lima2021jstransplant,zhong2025sigmodreuse,qin2019testmig,mobile_migration_combo_2021_2024, gao2024mut, compiler_testing_combo_2011_2014_2021, transplantation_clones_combo_2015_2017}. \tname{} applies this approach to WebAssembly by reusing regression tests from runtime repositories, confirming each test on the originally reported version, and then executing the same tests across independently implemented runtimes to check for divergences. In JavaScript engines, researchers transplant regression tests across engines with lightweight adapters, normalized flags and outputs \cite{lima2021jstransplant}. In database systems, suites are unified across engines by converting formats, standardizing configuration, and factoring out engine-specific behavior so transplanted tests execute across SQLite, PostgreSQL, and DuckDB \cite{zhong2025sigmodreuse}. Similarly, \tname{} invokes common entry points across runtimes and normalizes outputs while mapping error messages to canonical categories, enabling comparable cross-runtime execution results. In mobile apps, migration techniques map recorded UI events and assertions between iOS and Android by adapting locators, widgets, and flows \cite{qin2019testmig}, \cite{mobile_migration_combo_2021_2024}. Human-in-the-loop methods migrate unit tests across languages and platforms by combining code mapping with test adaptation while maintaining test intent \cite{gao2024mut}. At the code-clone level, transplantation moves a donor clone into a recipient and reuses the donor’s tests to exercise the transplanted functionality in its new context \cite{transplantation_clones_combo_2015_2017}. Although compiler testing often relies on randomized differential testing, equivalence modulo inputs, and modern program generators, our work differs by starting from real regression tests rather than random program generation \cite{compiler_testing_combo_2011_2014_2021}.

\end{comment}




\begin{comment}
    
\vspace{0.75ex}
\noindent\textbf{Regression Tests for Transplantation.}~
“Test transplantation” reuses existing, fault-revealing tests across systems by adapting harnesses, aligning configuration/flags, and preserving (or adjusting) the oracle. Prior work demonstrates effectiveness across domains: JavaScript engines use lightweight adapters to normalize flags and outputs across engines~\cite{lima2021jstransplant}; database systems convert formats and standardize configuration so suites run across SQLite, PostgreSQL, and DuckDB~\cite{zhong2025sigmodreuse}; mobile-app studies map recorded UI events and assertions across iOS/Android~\cite{qin2019testmig,mobile_migration_combo_2021_2024}; human-in-the-loop methods migrate unit tests across languages/platforms while maintaining intent~\cite{gao2024mut}; and clone-level transplantation reuses donor tests to validate moved functionality~\cite{transplantation_clones_combo_2015_2017}. \tname{} applies this strategy to WebAssembly by (1) sourcing regression tests from runtime repositories, (2) confirming each test on the originally reported version, and (3) executing the same tests across independently implemented runtimes with normalized outputs and canonical error categories, revealing cross-runtime divergences. Unlike compiler-testing pipelines that emphasize random program generation (\eg, EMI/CSmith-style approaches)~\cite{compiler_testing_combo_2011_2014_2021}, \tname{} centers on \emph{real regression tests} as transplantable artifacts.

\vspace{0.5ex}
\noindent\textbf{Regression Tests for Fuzzing.}~
Regression tests also serve as \emph{seed corpora} that steer fuzzing toward bug-revealing regions of the space. Foundational datasets (Defects4J, BugsInPy, Bugs.jar, BEARS) provide reproducible fail–pass pairs and oracles that can bootstrap fuzz campaigns~\cite{just2014defects4j,widyasari2020bugsinpy,saha2018bugsjar,madeiral2019bears}. CI-mined BugSwarm highlights the scale and pitfalls (environment drift, flaky builds) of harvesting real regressions~\cite{tomassi2019bugswarm}; \tname{} addresses these risks by re-running fault-revealing tests on cited versions before use. Empirical studies show initial seeds materially affect outcomes and advocate publishing exact seeds and comparing across corpora~\cite{herrera2021seed,chen2019enfuzz}; \tname{} follows this guidance and will release its corpus. Corpus-shaping techniques—MoonShine’s distillation and FairFuzz’s coverage-guided prioritization~\cite{pailoor2018moonshine,lemieux2018fairfuzz,manes2021survey}—reduce redundancy and focus effort; \tname{} complements these by deduplicating while preserving failing inputs. In Wasm specifically, our approach augments mutation-based fuzzers~\cite{zhao2024wapplique,arteaga2022wasm,jiang2022wasmfuzzer,qian2025freewavm,zhang2025waltzz} with \emph{regression-test–derived seeds}, linking practical bug histories to systematic exploration and enabling downstream differential checks (Wadiff, Wasmaker)~\cite{zhou2023wadiff,cao2024wasmaker}.

\vspace{0.5ex}
\noindent\textbf{Transplantation of Tests.}~
(Integrated above as “Regression Tests for Transplantation.” We retain this short pointer to maintain continuity with prior drafts.)
\end{comment}

\noindent\textbf{Regression Tests for Transplantation.}~
\deleted[id=giorgi]{Test transplantation reuses existing tests across systems by adapting harnesses, configurations, and oracles. Prior studies demonstrate its effectiveness for JavaScript engines, databases, and mobile applications~\cite{lima2021jstransplant,zhong2025sigmodreuse,qin2019testmig,talebipour2021ui,zhang2024synthesis,gao2024mut,yang2011finding,le2014compiler,livinskii2020random,barr2015automated,zhang2017automated}. Prior work migrated regression tests between engines or platforms through lightweight adapters and standardized outputs. \tname{} applies the same concept to WebAssembly by reusing regression tests from runtime repositories, confirming them on their original versions, and executing them across independent runtimes to detect divergences. This bridges the idea of test reuse across implementations while maintaining test intent.}
\added[id=giorgi]{Test transplantation reuses existing tests across systems by adapting harnesses, configurations, and oracles, and has proven effective for JavaScript engines, databases, and mobile applications~\cite{lima2021jstransplant,zhong2025sigmodreuse,qin2019testmig,talebipour2021ui,zhang2024synthesis,gao2024mut,yang2011finding,le2014compiler,livinskii2020random,barr2015automated,zhang2017automated}. In JavaScript engines, regression tests are transplanted across engines to find new bugs~\cite{lima2021jstransplant}, and in database systems, a single test suite is reused across SQLite, PostgreSQL, and DuckDB~\cite{zhong2025sigmodreuse}. \tname{} applies this idea to \wasm{} by reusing regression tests from runtime repositories, confirming them on their original versions, and executing them across independent runtimes to detect divergences. Unlike these settings, \wasm{} is portable by design, so \tname{} runs the same test cases on every runtime without any adaptation.}








\vspace{0.75ex}
\noindent\textbf{Regression Tests for Fuzzing.}~
\deleted[id=giorgi]{Regression tests can also serve as valuable fuzzing seeds. Mutational fuzzers such as \wapplique{}, \wasmaker{}, \wadiff{}, \lwdiff{} and others transform existing modules to explore differential or buggy runtime behavior~\cite{zhao2024wapplique,cao2024wasmaker,zhou2023wadiff,zhou2025lwdiff,qian2025freewavm,zhang2025waltzz}. \tname{} builds on these by deriving fuzzing seeds from real-world regression tests and bug reports, linking historical failures with systematic fuzzing and multi-runtime validation.}
\added[id=giorgi]{Regression tests can also serve as fuzzing seeds. Mutational fuzzers such as \wapplique{}, \wasmaker{}, \wadiff{}, and others~\cite{zhao2024wapplique,cao2024wasmaker,zhou2023wadiff,zhou2025lwdiff,jiang2022wasmfuzzer,qian2025freewavm,zhang2025waltzz} transform existing modules to explore differential or buggy runtime behavior. \tname{} builds on these by changing only the seed source, deriving seeds from real regression tests and bug reports. This exposes how much existing corpora have saturated: \wapplique{} and \wasmaker{} originally reported \WappliquePaperBugCount{} and \WasmakerPaperBugCount{} bugs using \llvmbench{} and \wasmbench{}~\cite{zhao2024wapplique,cao2024wasmaker}, yet the same fuzzers on the same corpora now reveal far fewer new bugs, consistent with the saturation effect as runtimes mature~\cite{groce2020saturation,liyanage2023reachable}. Using \tname{} instead reveals the most new bugs of any corpus we evaluate. \deleted[id=giorgi]{ This aligns with seed selection studies showing that the choice of initial corpus strongly affects fuzzing outcomes~\cite{rebert2014optimizing,herrera2021seed}.}}
\added[id=ishrak]{Consistent with prior seed selection studies~\cite{rebert2014optimizing,herrera2021seed}, we also found that seed selection strongly influences fuzzing effectiveness. However, our study additionally shows that, in the WebAssembly domain, regression tests are an especially valuable seed source because they preserve proposal-specific semantic corner cases that continue to expose defects in independently and asynchronously developed runtimes.}



\vspace{0.75ex}
\noindent\added[id=giorgi]{\textbf{Differential Testing of WebAssembly Runtimes.}}~
\added[id=giorgi]{Differential testing runs the same module on several runtimes and treats output disagreements as bugs; \wadiff{}~\cite{zhou2023wadiff}, \wasmaker{}~\cite{cao2024wasmaker}, and \wapplique{}~\cite{zhao2024wapplique} follow this approach and serve as baseline fuzzers in our study. These prior work focus more on the fuzzing approach whereas we focus more on an empirical study of the seed corpora used by these fuzzers. Two recent techniques add LLM assistance: \lwdiff{}~\cite{zhou2025lwdiff} supports differential testing of WebAssembly runtimes, and DrWASI~\cite{zhang2025drwasi} targets WebAssembly System Interface (WASI) implementations. \lwdiff{} uses an LLM to generate mutators from the WebAssembly specification~\cite{zhou2025lwdiff}. DrWASI adapts behavior from corresponding POSIX APIs and generates WASM test cases that exclusively target only the WASI implementations~\cite{zhang2025drwasi}. Neither \lwdiff{} nor DrWASI reuses regression tests for fuzzing, whereas \tname{} proposes a new regression tests based seed corpus for fuzzing WebAssembly runtimes. The LLM-based fuzzers and \tname{} are complementary, and \tname{} can seed an LLM-driven generator like \lwdiff{}.}



\vspace{0.75ex}
\noindent\textbf{Other Related Works.}~
\deleted[id=giorgi]{Beyond regression-test reuse and fuzzing, several studies contribute to the broader understanding of WebAssembly reliability and fuzzing evaluation. Structure- and rule-guided Wasm generators such as WasmCFuzz, RGFuzz, and SwFuzz synthesize valid programs to exercise deep compiler and runtime behaviors~\cite{zhang2024wasmcfuzz,park2025rgfuzz,wang2024s}. Static and analysis-based approaches, including WasmChecker and WBSan, detect type and memory errors and undefined behaviors without relying on fuzzing~\cite{zhuang2024wasmchecker,wu2025wbsan}. Formal verification efforts such as WasmRef-Isabelle provide a mechanized reference interpreter for validation and reasoning~\cite{watt2023wasmref}. Broader surveys like the SoK ``Analysis Techniques for WebAssembly''~\cite{harnes2024sok} and studies on Random Testing of WASI~\cite{stanley2024random} contextualize the evolving Wasm testing landscape. Work on evaluation methodology also informs our design: comparative studies emphasize the importance of statistical rigor and reproducibility~\cite{klees2018evaluating,arcuri2011statpract,bohme2022reliability}, while shared benchmarking infrastructures such as FuzzBench and GreenBench~\cite{metzman2021fuzzbench,ounjai2023greenbench} demonstrate reproducible experimental practices. Mutation-analysis research highlights the importance of separating coverage and mutation phases~\cite{gorz2023mutation}, and recent studies on mutator-level operator design~\cite{kukucka2024mutator} and multi-dimensional fuzzing metrics~\cite{li2021unifuzz} offer guidance for controlled experimentation.}
\added[id=giorgi]{Other work targets WebAssembly reliability and fuzzing methodology more broadly. Structure and rule-guided generators such as WasmCFuzz, RGFuzz, and SwFuzz~\cite{zhang2024wasmcfuzz,park2025rgfuzz,wang2024s} build valid programs from a grammar or the specification. \tname{} instead starts from real regression tests that already concentrate in bug-prone code, so the two are complementary as starting corpora. WasmChecker~\cite{zhuang2024wasmchecker} detects type and memory errors through static analysis without fuzzing, while WBSan~\cite{wu2025wbsan} provides sanitization to support binary-only fuzzing. WasmRef-Isabelle~\cite{watt2023wasmref} provides a verified reference interpreter that can serve as a fuzzing oracle. The SoK on analysis techniques for WebAssembly~\cite{harnes2024sok} and work on random testing of WASI~\cite{stanley2024random} survey this space. Studies on statistical rigor and reliable benchmarking~\cite{klees2018evaluating,arcuri2011statpract,bohme2022reliability} and platforms such as FuzzBench and GreenBench~\cite{metzman2021fuzzbench,ounjai2023greenbench} establish good practice for fuzzing evaluation. We follow this guidance by running repeated, fixed-budget campaigns and releasing our seeds. Mutation-analysis and mutator-level studies~\cite{gorz2023mutation,kukucka2024mutator} further motivate our use of default mutator settings across trials.}